\chapter{General Theoretical Framework and Conventions}
\label{ch:intro}

% This chapter's Implementation-in-tpeanuts blocks use full descriptive
% paragraphs rather than the one-line module purpose used from Chapter 2
% onwards, so the shared \modulehead's blanket italics (meant for a short
% phrase) reads poorly here. Locally drop it to plain upright text -- italics
% stay available for deliberately highlighting a single term -- and restore
% the shared definition at the end of the chapter so later chapters are
% unaffected.
\let\savedmoduleheadchOne\modulehead
\renewcommand{\modulehead}[2]{%
  \noindent\pkgpath{#1}\par
  \noindent#2\par\medskip
}

\section{Purpose and how to read this document}

\texttt{tpeanuts} is a \texttt{PyTorch} rewrite (GPU-first, with batching
support) of the original \texttt{peanuts} code \parencite{LucentePeanuts},
extended with BSM scenarios (NSI and a 3+1 sterile neutrino), a proprietary
perturbative evolution scheme, and interface layers to external reference
engines (\texttt{nuSQuIDS}, \texttt{MCEq}, Honda/HKKM tables,
\texttt{pymsis}).

Every section in the chapters that follow corresponds to \emph{one source
file} in the repository and always follows the same structure:

\begin{itemize}
  \item \modulehead{package/module.py}{one line describing its purpose}
    identifies the exact file.
  \item The blue box \textbf{Theoretical Foundation} contains the physics and
    mathematical derivations that justify what the module computes.
  \item The green box \textbf{Implementation in \texttt{tpeanuts}} translates
    that theory into the actual functions/classes, with their essential
    signature and the sign, unit, or tensor-shape conventions that must be
    respected to use them correctly.
  \item The orange box \textbf{Approximations and Limitations} (when
    applicable) documents the range of validity, limiting cases, and
    numerical simplifications.
  \item The grey box \textbf{Key References} links the specific bibliography
    for that piece of physics.
  \item Central results are \keyresult{highlighted}; sections that do not yet
    have a complete derivation are explicitly flagged as pending and are
    listed in the pending-tasks index right after the table of contents.
\end{itemize}

\section{The physics problem: neutrino oscillation in matter}

A neutrino is produced and detected in a flavour eigenstate
$\ket{\nu_\alpha}$, $\alpha \in \{e,\mu,\tau\}$ (plus $\ket{\nu_s}$ in the
sterile sector, \cref{ch:core-bsm}), but propagates as a superposition of
mass eigenstates $\ket{\nu_i}$. The PMNS mixing matrix
\parencite{Pontecorvo1957,MakiNakagawaSakata1962} connects the two bases,
\begin{equation}
  \ket{\nu_\alpha} = \sum_i U_{\alpha i}\,\ket{\nu_i}.
\end{equation}
In vacuum each mass eigenstate evolves with a kinetic phase $e^{-iE_i L}$;
when crossing matter an additional effective potential from coherent forward
interaction appears (the MSW effect,
\parencite{Wolfenstein1978,MikheyevSmirnov1985}). In both cases the evolution
is summarised by an effective Hamiltonian $H(x)$ acting on the vector of
flavour amplitudes,
\begin{equation}
  i \frac{d}{dx}\, \psi(x) = H(x)\, \psi(x),
  \qquad
  \psi = \begin{pmatrix}\psi_e\\ \psi_\mu\\ \psi_\tau\end{pmatrix},
  \label{eq:schrodinger-flavour}
\end{equation}
where $x$ is the dimensionless propagation coordinate defined in
\cref{sec:units}. This document organises \emph{all} propagation code around
three objects that appear systematically in every module:

\begin{enumerate}
  \item the \textbf{Hamiltonian} $H = H_{\rm kin} + H_{\rm mat}$
    (\cref{ch:core-common,ch:core-bsm}),
  \item the \textbf{evolution operator} (evolutor) $S$ that formally solves
    \labelcref{eq:schrodinger-flavour},
    $\psi(x_2) = S(x_2, x_1)\,\psi(x_1)$, obtained either exactly (constant
    density, \cref{ch:core-perturbative}) or numerically via segmented
    exponentiation (\cref{ch:core-numerical}),
  \item the \textbf{transition probability}
    $P(\nu_\alpha \to \nu_\beta) = |S_{\beta\alpha}|^2$
    (\cref{ch:core-common}), the final observable quantity.
\end{enumerate}

The \emph{Physical Media} chapters (\cref{ch:vacuum,ch:solar,ch:earth,%
ch:atmosphere}) instantiate these three objects for each concrete
propagation environment (vacuum, solar interior, Earth, atmosphere), and the
\emph{External Engines} chapters document how those results are
cross-checked against independent reference codes.

\section{Notation and global conventions}
\label{sec:units}

\begin{theorybox}
Physics is worked out in natural units $\hbar = c = 1$ unless explicitly
converted to laboratory units. Public input quantities always use:
\begin{itemize}
  \item neutrino energy $E$ in $\mathrm{MeV}$,
  \item squared mass differences $\Delta m^2$ in $\mathrm{eV}^2$,
  \item path lengths $L$ in metres,
  \item electron density $n_e$ in $\mathrm{mol\,cm^{-3}}$ (molar electron
    density, not mass density).
\end{itemize}
The Hamiltonian is made dimensionless by introducing an \emph{evolution
scale} $L_{\rm scale}$ (by default the Earth radius $R_\oplus$) and the
propagation coordinate
\begin{equation}
  x \equiv \frac{L}{L_{\rm scale}},
\end{equation}
so that $H(x)$ in \labelcref{eq:schrodinger-flavour} is unitless. The two
Hamiltonian terms are made dimensionless as
\begin{align}
  k_i &= L_{\rm scale}\,\frac{\Delta m^2_{i1}}{2E}\,\frac{1}{\hbar c}
       &&\text{(kinetic term, see \cref{ch:core-common}),}\\
  V &= \pm\sqrt{2}\,G_F\, n_e\, L_{\rm scale}
       &&\text{(MSW matter potential, see \cref{ch:core-common}),}
\end{align}
with sign $+$ for neutrinos and $-$ for antineutrinos (the standard sign
convention for the charged-current term).
\end{theorybox}

\begin{implbox}
\begin{itemize}
  \item \modulehead{util/constant.py}{single source of truth for physical
    constants: $G_F$ in $\mathrm{MeV^{-2}}$, $\hbar$, $c$, $\hbar c$, $N_A$,
    Earth and Sun radii, the astronomical unit, the proton mass, and the
    mass-to-nucleon-density conversion factor.}
  \item \modulehead{util/context.py}{\pyclass{RuntimeContext} bundles
    \code{(device, dtype)} and is passed explicitly throughout the stack
    instead of duplicating those two arguments in every function.}
  \item \modulehead{util/type.py}{tensor conversion and validation: resolving
    the complex dtype paired with a real dtype
    (\pyfunc{cdtype\_from\_real}), building/validating 3- or 4-flavour state
    vectors (\pyfunc{state\_tensor}), and broadcasting flavour vectors to
    batched shapes (\pyfunc{broadcast\_last3}).}
  \item \modulehead{util/torch\_util.py}{default device/dtype resolution,
    generic tensor broadcasting, batch flattening, and the convention
    \pyfunc{enforce\_identity\_for\_zero\_length} (a zero-length evolutor
    must collapse exactly to the identity, which numerical exponentials do
    not automatically guarantee in the $x\to 0$ limit).}
  \item \modulehead{core/common/neutrino.py}{canonical flavour-index
    convention: $e\to0$, $\mu\to1$, $\tau\to2$ (plus $s\to3$ in the sterile
    sector), used throughout the code to index tensors.}
\end{itemize}
\end{implbox}

\subsection*{Probability matrix convention}

\begin{theorybox}
Given the flavour evolutor $S$, the final amplitude is
$\psi_\beta(x_2) = \sum_\alpha S_{\beta\alpha}\,\psi_\alpha(x_1)$, so
\begin{equation}
  \keyresult{P(\nu_\alpha \to \nu_\beta) = \left|S_{\beta\alpha}\right|^2},
\end{equation}
with the \emph{initial flavour} index in the column and the \emph{final
flavour} index in the row. This is the convention fixed, once and for the
whole project, by \texttt{core.common.probability}.
\end{theorybox}

\begin{implbox}
\modulehead{core/common/probability.py}{full derivation in
\cref{ch:core-common}. Every \code{*\_probability\_state(...)} function in
the medium modules (\code{vacuum\_probability\_state},
\code{solar\_probability\_state}, \code{earth\_probability\_state},
\code{atmosphere\_probability\_state}) delegates to the utilities in this
module to keep the index convention identical across the whole project.}
\end{implbox}

\section{Coherent states versus incoherent mixtures}
\label{sec:coherent-vs-incoherent}

\begin{theorybox}
The code systematically distinguishes two representations of the state being
propagated, selected with the \code{massbasis} parameter:
\begin{description}
  \item[Coherent propagation] (\code{massbasis=False}). The state is a
    flavour-amplitude vector $\psi_\alpha$ carrying physical relative
    phases. The evolutor is applied directly,
    $\psi_{\rm final} = S\,\psi_{\rm inicial}$, and
    $P_\alpha = |\psi_{{\rm final},\alpha}|^2$. This is the correct regime
    when trajectories are short compared to the coherence length (vacuum,
    atmosphere, Earth in a single leg).
  \item[Incoherent mass mixture] (\code{massbasis=True}). The state is a
    vector of real weights $w_i \geq 0$, $\sum_i w_i = 1$, over mass
    eigenstates \emph{without} relative-phase information. The final flavour
    probability is
    \begin{equation}
      P_\alpha = \sum_i \left| \bigl(S\, U_{\rm PMNS}\bigr)_{\alpha i}\right|^2 w_i.
    \end{equation}
    This is the physically correct regime for solar neutrinos: the spread at
    the production point and the enormous Sun--Earth distance completely
    decohere the mass superposition before it reaches Earth
    (\cref{ch:solar}), so only the populations $w_i$ survive, not the
    phases.
\end{description}
\end{theorybox}

\begin{notebox}
The choice \code{massbasis=True/False} is not a numerical detail: it is a
physical statement about whether the neutrino wave-packet's coherence length
is larger or smaller than the scale of the problem. Using the incoherent
mode on a short leg (e.g.\ only inside the atmosphere) would discard real
interference; using the coherent mode for Sun--Earth propagation would
overstate coherence and produce spurious interference patterns that are not
observed.
\end{notebox}

\section{Named presets and parameter traceability}

\begin{implbox}
\modulehead{core/common/presets.py}{single registry
(\pyfunc{register\_preset}/\pyfunc{get\_preset}/\pyfunc{list\_presets}) for
every named parameter set in the project: \pyclass{OSCILLATION\_PRESETS}
(mixing angles and $\Delta m^2$, for both the standard 3-flavour sector and
the 3+1 sterile extension) and \pyclass{NSI\_PRESETS} (non-standard
couplings). Each preset carries its physical justification and bibliographic
citation (e.g.\ \cite{Esteban2020NuFit} for standard best-fit oscillation
parameters, \cite{Esteban2018LMADark} for the \emph{LMA-Dark} preset, and
\cite{IceCube2022NSI} for the IceCube DeepCore NSI bounds), so any numerical
result can be traced back to its origin in the literature.}
\end{implbox}

\begin{notebox}
Counting CP phases for the sterile extension requires care: an $N$-generation
Dirac mixing matrix has $(N-1)(N-2)/2$ independent physical phases. For
$N=4$ (the 3+1 case) this gives exactly \keyresult{three} Dirac phases
($\delta_{13}$, $\delta_{14}$, $\delta_{24}$); the rotation $R_{34}$ is
always real because the fourth phase can be absorbed into a redefinition of
the charged-lepton fields, and hence sterile presets do not include
\code{delta34\_deg} (see \cref{ch:core-bsm}).
\end{notebox}

% Restore the shared \modulehead (one-line italic description) for every
% chapter from here on.
\let\modulehead\savedmoduleheadchOne
